\chapter{Coroutines I --- co\_await, co\_yield, co\_return, and the Promise Protocol}
\label{ch:c20-coroutines-i}

\begin{quote}
\emph{Coroutines are the third pillar and the most unusual: C++20 ships the low-level language machinery --- three keywords and a customization protocol --- but almost none of the high-level types you would actually use (\texttt{std::generator} and a \texttt{task} type are C++23 and beyond). This chapter explains the keywords, the compiler-generated state machine and coroutine frame, and the promise-type protocol.}
\end{quote}

A coroutine can \textbf{suspend}, hand control back to its caller, and later be \textbf{resumed} with its locals intact. C++20 gives you the engine, not the car: to use coroutines you implement a \textbf{promise type}.

\section{What a Coroutine Is}
\label{sec:c20-what-coroutine}

C++20 coroutines are \textbf{stackless}: the compiler transforms the function into a state machine whose state lives in a heap-allocated \textbf{coroutine frame}. Two dominant uses: \textbf{generators} (\texttt{co\_yield} a sequence on demand) and \textbf{tasks/async} (\texttt{co\_await} operations without blocking a thread).

\section{The Three Keywords}
\label{sec:c20-three-keywords}

\begin{center}
\begin{tabular}{ll}
\hline
\textbf{Keyword} & \textbf{Meaning} \\
\hline
\texttt{co\_await expr} & suspend until \texttt{expr} (awaitable) is ready \\
\texttt{co\_yield expr} & suspend and produce a value (sugar for \texttt{co\_await promise.yield\_value}) \\
\texttt{co\_return expr} & finish, delivering a final value \\
\hline
\end{tabular}
\end{center}

\begin{lstlisting}[language=C++,caption={Each keyword makes the enclosing function a coroutine}]
#include <coroutine>

Generator<int> counter();         // uses co_yield
Task<int>      fetch();           // uses co_await/co_return
Task<void>     fire_and_forget(); // uses co_return;
\end{lstlisting}

A coroutine cannot use plain \texttt{return}, be \texttt{constexpr}/\texttt{consteval}, be \texttt{main}, or be variadic.

\section{The Coroutine Frame and the Generated State Machine}
\label{sec:c20-frame}

On call, the compiler allocates the frame (heap, possibly elided via HALO), copies parameters by value, constructs the promise, calls \texttt{get\_return\_object()}, then evaluates \texttt{co\_await promise.initial\_suspend()}. A \textbf{\texttt{std::coroutine\_handle<Promise>}} resumes (\texttt{.resume()}), queries (\texttt{.done()}), reaches the promise (\texttt{.promise()}), and destroys (\texttt{.destroy()}).

\begin{lstlisting}[language=C++,caption={coroutine\_handle drives a coroutine from non-coroutine code}]
#include <coroutine>

template<typename Promise>
void drive(std::coroutine_handle<Promise> h) {
    while (!h.done()) h.resume();
    h.destroy();
}
\end{lstlisting}

\section{The Promise Type Protocol}
\label{sec:c20-promise-protocol}

The promise type is found via \texttt{std::coroutine\_traits<ReturnType, Args...>::promise\_type} (default: nested \texttt{ReturnType::promise\_type}).

\begin{center}
\begin{tabular}{ll}
\hline
\textbf{Promise member} & \textbf{Purpose} \\
\hline
\texttt{get\_return\_object()} & object returned to caller \\
\texttt{initial\_suspend()} & suspend or run immediately \\
\texttt{final\_suspend() noexcept} & keep frame alive or end \\
\texttt{unhandled\_exception()} & handle escaping exception \\
\texttt{return\_value(v)} / \texttt{return\_void()} & store result (exactly one) \\
\texttt{yield\_value(v)} & transport yielded value \\
\hline
\end{tabular}
\end{center}

\section{The Coroutine Lifecycle, Step by Step}
\label{sec:c20-lifecycle}

\begin{enumerate}
\item Call $\rightarrow$ frame allocated, parameters copied, promise constructed.
\item \texttt{get\_return\_object()} $\rightarrow$ caller's handle/value.
\item \texttt{co\_await initial\_suspend()} $\rightarrow$ suspend or run.
\item Body executes to a \texttt{co\_await}/\texttt{co\_yield} or completion.
\item \texttt{return\_value}/\texttt{return\_void} on completion.
\item \texttt{unhandled\_exception()} if the body throws.
\item \texttt{co\_await final\_suspend()} $\rightarrow$ usually \texttt{suspend\_always}.
\end{enumerate}

\section{initial\_suspend and final\_suspend}
\label{sec:c20-suspend-points}

\begin{lstlisting}[language=C++,caption={The two standard trivial awaitables drive start/end policy}]
#include <coroutine>

std::suspend_always initial_suspend() noexcept { return {}; }  // lazy (generator)
std::suspend_never  initial_suspend() noexcept { return {}; }  // eager (task)
std::suspend_always final_suspend()   noexcept { return {}; }  // keep frame alive
\end{lstlisting}

\texttt{initial\_suspend} defines lazy vs eager; \texttt{final\_suspend} must be \texttt{noexcept}.

\section{A Complete Minimal Generator}
\label{sec:c20-minimal-generator}

\begin{lstlisting}[language=C++,caption={A complete hand-written generator coroutine in C++20}]
#include <coroutine>
#include <iostream>
#include <utility>

template<typename T>
struct Generator {
    struct promise_type;
    using handle_type = std::coroutine_handle<promise_type>;

    struct promise_type {
        T current_value;
        Generator get_return_object() {
            return Generator{ handle_type::from_promise(*this) };
        }
        std::suspend_always initial_suspend() noexcept { return {}; }
        std::suspend_always final_suspend() noexcept { return {}; }
        std::suspend_always yield_value(T value) {
            current_value = std::move(value);
            return {};
        }
        void return_void() {}
        void unhandled_exception() { std::terminate(); }
    };

    handle_type h;
    explicit Generator(handle_type handle) : h(handle) {}
    Generator(const Generator&) = delete;
    Generator(Generator&& o) noexcept : h(std::exchange(o.h, {})) {}
    ~Generator() { if (h) h.destroy(); }

    bool next() { h.resume(); return !h.done(); }
    const T& value() const { return h.promise().current_value; }
};

Generator<int> sequence(int start, int end) {
    for (int i = start; i < end; ++i)
        co_yield i;
}

int main() {
    auto gen = sequence(0, 5);
    while (gen.next())
        std::cout << gen.value() << ' ';   // 0 1 2 3 4
}
\end{lstlisting}

\section{Exceptions and unhandled\_exception}
\label{sec:c20-coro-exceptions}

\begin{lstlisting}[language=C++,caption={Capturing the exception to rethrow on the consumer side}]
#include <coroutine>
#include <exception>

struct promise_type_with_errors {
    std::exception_ptr error;
    void unhandled_exception() noexcept {
        error = std::current_exception();
    }
    // consumer: if (h.promise().error) std::rethrow_exception(h.promise().error);
};
\end{lstlisting}

Because \texttt{final\_suspend} is \texttt{noexcept}, the frame is intact when the consumer inspects the captured exception.

\section{Professional Insights}
\label{sec:c20-coroutines-i-insights}

\textbf{C++20 ships the engine, not the car.} No standard \texttt{generator}/\texttt{task}/executor --- hand-roll or import (cppcoro, libcoro, Asio).

\textbf{Never capture a reference parameter across a suspension.} Parameters are copied into the frame, but a reference copies the reference; a temporary referent dangles after the first suspension. Pass by value.

\textbf{Make \texttt{final\_suspend} \texttt{noexcept} and usually \texttt{suspend\_always}.} It keeps the frame alive so the consumer can read the result or exception.

\textbf{Budget for the frame allocation in hot paths.} HALO may elide it, but not across TUs or type-erased handles; supply a custom promise \texttt{operator new} (pool/arena) and profile.

\textbf{Capture exceptions in the promise.} Store \texttt{std::current\_exception()} and rethrow from the consumer accessor rather than calling \texttt{std::terminate}.
